% Figure 4 — four conjunctive proofs on one release rail. Every stage uses a
% distinct glyph; failure is a visible break, not another prose row.
\begin{figure}[H]
\centering
\begin{tikzpicture}[pd figure,x=1cm,y=1cm]
  \draw[pd hairline,line width=1.15pt] (.65,1.78)--(14.55,1.78);
  \draw[pd focus arrow,line width=1.15pt] (14.55,1.78)--(15.25,1.78);
  \node[pd direct label,text=hhteal,anchor=west,font=\footnotesize\bfseries]
    at (15.02,2.04) {release};

  \foreach \x/\n in {2.0/1,5.55/2,9.10/3,12.65/4}{
    \node[font=\scriptsize\bfseries,text=hhgray] at (\x,3.82) {0\n};
    \draw[pd focus rule,line width=1.25pt] (\x,1.48)--(\x,2.08);
    \node[pd focus datum] at (\x,1.78) {};
  }

  \draw[pd rule] (1.35,2.62) rectangle (2.65,3.35);
  \draw[pd focus rule] (1.72,2.99) circle (.18);
  \draw[pd focus rule] (1.90,2.99)--(2.31,2.99);
  \node[pd panel title,anchor=north,align=center] at (2.0,1.30)
    {irreversible?};
  \node[pd axis label,anchor=north,align=center] at (2.0,.86)
    {artifact one key away};

  \draw[pd rule] (4.85,2.60) rectangle (5.35,3.36);
  \foreach \y in {2.76,2.98,3.20}{\draw[pd hairline] (4.95,\y)--(5.25,\y);}
  \draw[pd focus rule] (5.73,2.62).. controls (5.93,2.82) and (5.93,3.16)..(6.20,3.34);
  \draw[pd focus rule] (5.73,2.98).. controls (5.95,3.12) and (6.02,2.72)..(6.20,2.62);
  \node[pd panel title,anchor=north,align=center] at (5.55,1.30)
    {state or m\^etis?};
  \node[pd axis label,anchor=north,align=center] at (5.55,.86)
    {flatten state; link trace};

  \draw[pd rule] (8.45,2.74)--(9.02,3.32)--(9.74,2.62);
  \draw[pd focus rule,line width=1.2pt] (8.62,2.87)--(8.92,2.65)--(9.54,3.22);
  \node[pd panel title,anchor=north,align=center] at (9.10,1.30)
    {could it lie?};
  \node[pd axis label,anchor=north,align=center] at (9.10,.86)
    {resolve beyond self-report};

  \draw[pd rule] (12.05,2.62) arc[start angle=210,end angle=-30,x radius=.70,y radius=.44];
  \draw[pd focus arrow] (12.22,2.78) arc[start angle=210,end angle=40,x radius=.52,y radius=.30];
  \node[font=\scriptsize\bfseries] at (12.65,2.99) {why};
  \node[pd panel title,anchor=north,align=center] at (12.65,1.30)
    {authority legible?};
  \node[pd axis label,anchor=north,align=center] at (12.65,.86)
    {return a named reason};

  \draw[pd caution rule] (5.95,.25)--(7.05,.25);
  \draw[pd caution rule] (7.55,.25)--(8.65,.25);
  \draw[pd caution rule,line width=1.2pt] (7.13,.10)--(7.47,.40);
  \draw[pd caution rule,line width=1.2pt] (7.47,.10)--(7.13,.40);
  \node[pd note,text=hhamber,anchor=north] at (7.30,.02)
    {one missing proof holds the surface};
\end{tikzpicture}
\caption{A read-surface earns release through four independent proofs, in
order: irreversible action exposes its artifact; administrative state remains
separate from verbatim m\^etis; a claim resolves beyond the author’s
self-report; and every coordinator act returns a named reason. The rail is
conjunctive: one missing proof breaks release.}
\label{fig:four-questions}
\end{figure}
